% =========================================================================
% Capítulo 6 — Análisis de alternativas
% Método: para cada ámbito, qué se necesita, qué opciones había, cuál se
% elige y por qué. El punto de partida es siempre la arquitectura de
% TrackBites (proyecto previo del autor); las decisiones aún abiertas se
% declaran como tales, con sus candidatas y el criterio que decidirá.
% Fuentes: propuesta original (§4), dev-journal, CLAUDE.md y el propio
% estado del repositorio. La solución resultante se describe en el cap. 8.
% =========================================================================

\chapter{Análisis de alternativas}
\label{ch:alternativas}

Este capítulo recorre las decisiones tecnológicas del proyecto: para cada ámbito se expone qué se necesita, qué opciones se han considerado, cuál se elige y por qué.
El análisis no parte de cero.
El autor desarrolló previamente TrackBites \parencite{deleanu-trackbites} (apartado \ref{sec:contexto-marco}), y aquella arquitectura funciona aquí como línea base: cada elección se plantea como «lo ya dominado» frente a lo que este producto exige, y solo se salta donde el salto aporta.
La propuesta de este \ac{TFG} concentró ese salto en cuatro mejoras: TypeScript en todo el código, lógica compartida entre web y móvil, una base de datos relacional con capacidades geoespaciales y vectoriales, y el cierre del ciclo de entrega con despliegue automático.

Tres criterios transversales, fijados en capítulos anteriores, atraviesan todas las decisiones.
Primero, el \ac{TFG} debe demostrar las competencias del grado, así que las piezas centrales se construyen a mano en lugar de delegarse en plataformas que las regalan hechas (sección \ref{sec:alt-supabase}).
Segundo, la restricción de coste cero: solo se adoptan servicios con capa gratuita suficiente para un uso académico (riesgo R9, capítulo \ref{ch:riesgos}).
Tercero, el \ac{RGPD} condiciona dónde viven los datos (región europea) y qué servicios pueden tratarlos (riesgo R7).
Las decisiones diferidas se recogen con sus candidatas y con el criterio que las cierra.
La solución resultante de este análisis se describe en el capítulo \ref{ch:solucion}.

\section{Lenguaje: TypeScript en todo el código}
\label{sec:alt-lenguaje}

TrackBites se escribió en JavaScript.
Aquí todo el código —servidor, web, móvil y paquetes compartidos— pasa a TypeScript, la decisión que más mantenibilidad aporta a medio y largo plazo: el compilador detecta antes de ejecutar errores que en JavaScript solo aparecen en ejecución \parencite{typescript-docs}, y los tipos son la base técnica de la validación compartida de la sección \ref{sec:alt-validacion}.
La configuración es deliberadamente estricta (modo \eng{strict} con comprobaciones adicionales) y única para todo el repositorio, de modo que ningún paquete pueda relajarla por su cuenta.

\section{Estructura del repositorio: \eng{monorepo} con pnpm}
\label{sec:alt-monorepo}

Compartir tipos, validación y cliente de \ac{API} entre tres aplicaciones exige que todo el código viva junto: un único repositorio con las aplicaciones (\texttt{apps/}) y los paquetes compartidos (\texttt{packages/}), el llamado \eng{monorepo}.
La alternativa de un repositorio por aplicación se descarta porque cada pieza compartida tendría que publicarse y versionarse para poder consumirse, un peaje continuo sin sentido para una sola persona.

Dentro del \eng{monorepo} compitieron dos herramientas.
La propuesta original apuntaba a Turborepo, un orquestador que acelera \eng{monorepos} grandes con caché de tareas y compilaciones incrementales \parencite{turborepo-docs}.
Al montar el repositorio se optó en su lugar por los \eng{workspaces} de pnpm sin orquestador \parencite{pnpm-docs}.
La arquitectura elegida explica por sí sola esa renuncia: los paquetes compartidos se consumen \emph{sin compilar} (sus puntos de entrada apuntan directamente al código fuente y cada aplicación los compila con su propia herramienta), así que apenas quedan tareas que orquestar ni compilaciones que cachear, y la capa extra de Turborepo pagaría complejidad sin devolver velocidad en un equipo de una persona.
pnpm aporta además la reproducibilidad que el proyecto se exige: versiones de dependencias fijadas exactas (sin rangos) con el archivo de bloqueo versionado en el repositorio, y versiones de Node y del propio gestor clavadas en la configuración, de modo que cualquier máquina reconstruye el mismo entorno.
Queda preparada una contingencia documentada: si el empaquetador de React Native no tolerase los enlaces simbólicos con los que pnpm instala, una línea comentada en la configuración cambia a una instalación plana equivalente a la de npm sin perder el almacén rápido de pnpm.

\section{Base de datos: PostgreSQL con PostGIS y pgvector}
\label{sec:alt-bd}

Es la decisión más relevante del conjunto, y la única que hereda una lección previa documentada.
TrackBites usó MongoDB con Mongoose, y su memoria dejó registrado un compromiso consciente: aquel dominio, estrictamente jerárquico, encajaba mejor en una base relacional, y la integridad referencial y los borrados en cascada hubo que emularlos a mano con ganchos del propio Mongoose.
\bukit lleva esa conclusión a la práctica.
Su dominio es un grafo social —personas que se siguen, listas, metas, publicaciones ligadas a listas— con dos necesidades especiales: buscar personas por cercanía geográfica y por afinidad de intereses.
Un modelo relacional resuelve el grafo con claves foráneas e integridad nativas, y aprovecha la base sólida en \ac{SQL} del autor.

Las dos necesidades especiales no obligan a añadir motores dedicados, porque PostgreSQL las cubre con extensiones: PostGIS añade tipos y consultas geoespaciales —búsqueda por radio con índice espacial— \parencite{postgis-docs}, y pgvector añade columnas vectoriales y búsqueda por similitud \parencite{pgvector-docs} sobre las representaciones numéricas (\eng{embeddings}) de metas e intereses con las que se calcula la afinidad.
La alternativa de sumar servicios especializados (una base vectorial y un servicio geoespacial junto a la base principal) multiplicaría las piezas que operar, pagar y sincronizar; un único PostgreSQL cubre las tres necesidades a la vez.

Para producción se elige Neon, un PostgreSQL gestionado \parencite{neon-docs}, entre las candidatas de la propuesta (Neon, Supabase, Railway): la residencia europea del dato es el criterio registrado de esa elección, y su capa gratuita satisface además la restricción de coste cero.
La disposición operativa de la base —Docker en desarrollo, pruebas e integración continua; Neon en producción— se describe en el capítulo \ref{ch:solucion}.

\section{Acceso a datos: Prisma}
\label{sec:alt-orm}

Como capa de acceso a datos (\ac{ORM}) se adopta Prisma, que aporta migraciones versionadas del esquema y generación automática de los tipos TypeScript de cada entidad \parencite{prisma-docs}: cada cambio del modelo queda versionado junto al código, y los tipos del servidor no pueden desincronizarse de la base.
Sustituye a Mongoose, la capa equivalente del mundo MongoDB, que queda descartada junto con su motor.
Prisma no cubre las consultas geoespaciales ni vectoriales de las extensiones; esas consultas concretas se escriben en \ac{SQL} crudo desde el propio Prisma, una convivencia prevista desde el diseño.
Las extensiones se declaran además en el propio esquema de Prisma, de modo que el código y la base nunca discrepen sobre qué hay activado.

\section{Servidor de la \acs{API}: Express 5 frente a NestJS}
\label{sec:alt-servidor}

La propuesta dejó abierta la elección entre NestJS y «Express bien estructurado», y el arranque del repositorio la cerró a favor de Express en su versión 5.
NestJS ofrece una arquitectura completa y opinada —módulos, inyección de dependencias, decoradores— \parencite{nestjs-docs}; Express es minimalista y deja la estructura en manos del proyecto \parencite{express-docs}.
Esa diferencia basta para sostener la elección.
El valor académico de este \eng{backend} está justamente en lo hecho a mano —autenticación con \ac{JWT} y rotación de los \eng{tokens} de refresco, doble capa de limitación de peticiones, manejo de errores centralizado: patrones que el autor ya domina de TrackBites, precisamente sobre Express—, y NestJS los entregaría envueltos en las abstracciones del marco, restando demostración y sumando su propia curva de aprendizaje a la del resto de la pila (riesgo R6, capítulo \ref{ch:riesgos}).
\section{Pruebas de la \acs{API} y registro de sucesos}
\label{sec:alt-pruebas-registro}

Dos decisiones de menor calado acompañan al servidor, ambas con alternativa heredada de TrackBites o de su ecosistema.
Las pruebas de integración de la \ac{API} se escriben con Vitest \parencite{vitest-docs} y Supertest, frente al precedente de TrackBites (Jest con Supertest).
Dos razones cierran la elección: un único ejecutor de pruebas para todo el \eng{monorepo} (los paquetes compartidos ya prueban con Vitest) y el hecho de que la \ac{API} se escribe como módulos ECMAScript, un terreno donde Jest exige una configuración frágil.
El registro de sucesos se resuelve con pino \parencite{pino-docs} y su \eng{middleware} pino-http, frente a winston: pino produce de forma nativa registro estructurado en \ac{JSON} con niveles de gravedad, su \eng{middleware} \ac{HTTP} asigna un identificador a cada petición para seguirla en los registros, y su rendimiento es mejor.

\section{La alternativa global: Supabase}
\label{sec:alt-supabase}

Antes de elegir pieza a pieza cabía una enmienda a la totalidad.
Supabase empaqueta en un solo producto casi todo lo que este \eng{backend} construye: PostgreSQL gestionado con PostGIS y pgvector, autenticación, almacenamiento de archivos, tiempo real y funciones en la nube \parencite{supabase-docs}.
Su ventaja es real y se reconoce: desarrollo mucho más rápido y menos infraestructura que mantener.
Se descarta por la razón inversa: menos \eng{backend} hecho a mano, que es justo lo que el autor domina y quiere demostrar.
Un \ac{TFG} cuyo \eng{backend} fuese configuración de Supabase acreditaría integración, no diseño; el \eng{backend} propio con Node, Prisma y PostgreSQL es el núcleo que demuestra las competencias del grado.
Supabase queda anotada como vía de aceleración legítima para una eventual segunda versión, ya fuera de este \ac{TFG}.

\section{Validación compartida: Zod frente a Joi}
\label{sec:alt-validacion}

TrackBites validaba con Joi solo en el servidor, sin validación compartida con el cliente.
Aquí la validación vive una única vez, en el paquete compartido \texttt{packages/core}, escrita con Zod: cada esquema valida los datos en ejecución y de él se infiere automáticamente el tipo TypeScript \parencite{zod-docs}, de modo que servidor, web y móvil comparten literalmente la misma regla y el mismo tipo.
Con Joi, validación y tipos habrían tenido que mantenerse por separado \parencite{joi-docs}: el doble mantenimiento que este diseño elimina es exactamente el porqué del cambio.
La regla de proyecto lo fija: ningún esquema se duplica en las aplicaciones.

\section{Cliente web: React y Vite, con TanStack Query}
\label{sec:alt-web}

La base del cliente web se conserva de TrackBites: React con Vite, ahora en TypeScript.
Cambiar de marco (por ejemplo, a uno con renderizado en el servidor) no aportaría a una aplicación que vive tras un inicio de sesión, y desperdiciaría la experiencia acumulada.
El cambio está en el estado de servidor.
En TrackBites la caché de datos remotos se gestionaba a mano con Redux Toolkit, podando el árbol de estado local para evitar recargas; funcionó, pero a costa de mucho código repetitivo de carga, error y caché.
Para un producto cuyo centro es un \eng{feed} con desplazamiento infinito, esa gestión manual se descarta en favor de TanStack Query, que resuelve de serie la caché, los reintentos, la invalidación y la paginación infinita del estado de servidor \parencite{tanstack-query-docs}.
La biblioteca de estado global de interfaz (Redux Toolkit con Reselect, en consideración por la experiencia previa) se decide al montar el cliente web, con el primer estado global real como caso de prueba.
Por el mismo principio, el andamiaje de las aplicaciones web y móvil se genera con sus herramientas oficiales en el momento de abordarlas, de modo que la configuración nace alineada con la versión vigente de cada herramienta.

\section{Cliente móvil: React Native con Expo}
\label{sec:alt-movil}

La aplicación móvil es objetivo secundario (capítulo \ref{ch:objetivos}), pero su elección condiciona la arquitectura desde el primer día.
Se elige React Native con Expo.
React Native permite escribir la interfaz móvil en el mismo lenguaje y paradigma que la web \parencite{react-native-docs}, con lo que los paquetes compartidos (tipos, validación, cliente de \ac{API}, lógica de estado) sirven a ambas plataformas sin traducción; el desarrollo nativo por plataforma —o un marco multiplataforma de otro lenguaje— rompería exactamente esa compartición, que es el principio rector de la arquitectura.
Expo añade la cadena de entrega: compilación en la nube, distribución de versiones de prueba y actualizaciones remotas sin pasar por las tiendas \parencite{expo-docs}.
La única fricción conocida (el empaquetador de React Native frente a los enlaces simbólicos de pnpm) tiene su contingencia preparada en la configuración del \eng{monorepo} (sección \ref{sec:alt-monorepo}).

\section{Servicios externos: enfoque decidido, proveedor abierto}
\label{sec:alt-servicios}

Varias funciones dependen de servicios de terceros.
En todas ellas el análisis distingue dos niveles: el \emph{enfoque}, decidido y condicionante del diseño, y el \emph{proveedor} concreto, que se elige en el momento de la integración con dos criterios fijos: capa gratuita suficiente allí donde exista (riesgo R9) y, cuando trate datos personales, garantías europeas de protección de datos (riesgo R7).
Cuando un servicio puede no ofrecer capa gratuita utilizable —el caso del vídeo y de la \ac{IA}—, esa dependencia de pago es justamente lo que relega su función a objetivo secundario (capítulo \ref{ch:objetivos}), con versiones simuladas como contingencia (riesgo R9).

\paragraph{Estimación de plazos y afinidad con \ac{IA}.}
El enfoque decidido es dirigido por eventos: el modelo de lenguaje (\ac{LLM}) se llama únicamente cuando la persona crea o edita una lista o pide una estimación; el sondeo periódico se descarta por costoso e innecesario.
La llamada se ejecuta como trabajo en segundo plano en una cola sobre Redis (BullMQ \parencite{bullmq-docs}) para no bloquear la petición, con respuesta estructurada en formato \ac{JSON} y con caché: si la lista no cambió, no se recalcula.
Las representaciones vectoriales de metas e intereses se generan al crearlos y se guardan en pgvector (sección \ref{sec:alt-bd}).
El proveedor del modelo queda abierto entre la \ac{API} de Anthropic y la de OpenAI; se elegirá al integrar la función comparando el coste por uso y las garantías europeas de protección de datos, los dos criterios fijados al inicio de esta sección.

\paragraph{Verificación de identidad.}
Construir un reconocimiento propio de documentos queda descartado de plano: se integra un proveedor especializado (candidatos: Stripe Identity, Veriff, Onfido) del que se guarda únicamente el estado de la verificación y una clave de referencia, nunca la imagen del documento, por minimización de datos (\ac{RGPD}).
La integración entra en el alcance en modo de pruebas; la contratación del servicio en producción queda fuera (capítulo \ref{ch:objetivos}).

\paragraph{Mensajería en tiempo real.}
El chat exige un canal bidireccional persistente entre cliente y servidor.
Las candidatas son Socket.IO \parencite{socketio-docs}, gestionado sobre el propio servidor, o un servicio equivalente ya gestionado; se decidirá al abordar la mensajería, con la capa gratuita como criterio.

\paragraph{Imágenes y vídeo.}
Las imágenes irán a un almacenamiento de objetos (Cloudflare R2 o S3, abierto).
El vídeo —objetivo secundario desde el 16/07/2026— irá a un servicio gestionado de vídeo (candidatos: Mux o Cloudflare Stream, abierto).

\paragraph{Notificaciones y anuncios.}
Los recordatorios de metas llegarán por notificaciones \eng{push} (Expo Push o Firebase Cloud Messaging, abierto).
La publicidad no se programa desde cero: el producto prevé una red de anuncios existente (AdMob en móvil; AdSense o Ad Manager en web, abierto) más una tabla propia de publicaciones patrocinadas que la ordenación del \eng{feed} inyecta.
Dentro de este \ac{TFG} solo se implementa la tabla propia de patrocinados: la integración de la red comercial queda fuera del alcance (capítulo \ref{ch:objetivos}).

\section{Entrega: de la integración continua heredada al despliegue continuo}
\label{sec:alt-entrega}

De TrackBites se hereda la cadena de comprobaciones automáticas (\eng{pipeline}) de integración continua en GitHub Actions, con estrategia de fallo rápido —análisis estático, comprobación de tipos y pruebas, con las pruebas de extremo a extremo de Playwright como última barrera—, cuya materialización concreta se describe en el capítulo \ref{ch:solucion}.
El salto nuevo es el despliegue continuo: en TrackBites se omitió deliberadamente, y aquí se completa el ciclo, de modo que todo cambio que pasa el \eng{pipeline} llega solo a producción.
El \eng{backend} se empaqueta con Docker, que sustituye el despliegue manual anterior por entornos reproducibles \parencite{docker-docs}.
Los destinos concretos se eligen al abordar el despliegue: la \ac{API}, entre Render, Railway y Fly.io; la web, entre Vercel y Cloudflare Pages; la aplicación móvil, por la cadena de Expo.
La base de datos es la excepción: su destino ya está decidido en Neon (sección \ref{sec:alt-bd}).
El criterio para cerrar los abiertos vuelve a ser el mismo: capa gratuita suficiente para un uso académico.

\section{Resumen del análisis}
\label{sec:alt-resumen}

La tabla \ref{tab:resumen-alternativas} condensa el capítulo: qué se elige, frente a qué alternativa y en qué estado queda cada decisión.

\begingroup
\footnotesize
\setlength{\tabcolsep}{4pt}
\begin{longtable}{@{}>{\raggedright\arraybackslash}p{31mm}>{\raggedright\arraybackslash}p{44mm}>{\raggedright\arraybackslash}p{44mm}>{\raggedright\arraybackslash}p{28mm}@{}}
  \caption{Resumen del análisis de alternativas.}
  \label{tab:resumen-alternativas}\\
  \toprule
  \textbf{Ámbito} & \textbf{Elección} & \textbf{Frente a} & \textbf{Estado} \\
  \midrule
  \endfirsthead
  \multicolumn{4}{c}{\emph{\tablename~\thetable{} — continuación de la página anterior}}\\
  \toprule
  \textbf{Ámbito} & \textbf{Elección} & \textbf{Frente a} & \textbf{Estado} \\
  \midrule
  \endhead
  \midrule
  \multicolumn{4}{r@{}}{\emph{Continúa en la página siguiente…}}\\
  \endfoot
  \bottomrule
  \endlastfoot
  Lenguaje & TypeScript transversal & JavaScript (TrackBites) & Decidida \\
  \addlinespace
  Repositorio & \eng{Monorepo} con \eng{workspaces} de pnpm & Turborepo; repositorios separados* & Decidida \\
  \addlinespace
  Base de datos & PostgreSQL + PostGIS + pgvector & MongoDB + Mongoose; motores especializados aparte & Decidida \\
  \addlinespace
  Base de datos en producción & Neon (región europea) & Supabase; Railway & Decidida \\
  \addlinespace
  Acceso a datos & Prisma (geo/vector en \acs{SQL} crudo) & Mongoose & Decidida \\
  \addlinespace
  Servidor de la \acs{API} & Express 5 estructurado a mano & NestJS & Decidida \\
  \addlinespace
  Pruebas de la \acs{API} & Vitest + Supertest & Jest + Supertest (TrackBites) & Decidida \\
  \addlinespace
  Registro de sucesos & pino + pino-http & winston & Decidida \\
  \addlinespace
  Plataforma \eng{backend} & \eng{Backend} propio (Node + Prisma + PostgreSQL) & Supabase como plataforma integral & Decidida \\
  \addlinespace
  Validación y tipos & Zod compartido en \texttt{packages/core} & Joi solo en el servidor & Decidida \\
  \addlinespace
  Cliente web & React + Vite + TanStack Query & Caché manual con Redux (TrackBites) & Decidida \\
  \addlinespace
  Estado global de interfaz & Pendiente (Redux Toolkit + Reselect en consideración) & — & Abierta \\
  \addlinespace
  Cliente móvil & React Native + Expo & Desarrollo nativo por plataforma* & Decidida \\
  \addlinespace
  \acs{IA} de plazos y afinidad & Por eventos + cola en segundo plano & Sondeo periódico & Enfoque decidido; proveedor abierto \\
  \addlinespace
  Verificación de identidad & Proveedor externo, en modo de pruebas & Reconocimiento propio & Enfoque decidido; proveedor abierto \\
  \addlinespace
  Mensajería en tiempo real & Canal bidireccional persistente & — & Enfoque decidido; proveedor abierto \\
  \addlinespace
  Imágenes y vídeo & Almacén de objetos + servicio de vídeo gestionado & — & Enfoque decidido; proveedor abierto \\
  \addlinespace
  Notificaciones \eng{push} & Expo Push o Firebase Cloud Messaging & — & Enfoque decidido; proveedor abierto \\
  \addlinespace
  Publicidad & Tabla propia de patrocinados (red existente prevista) & Publicidad programada desde cero & Enfoque decidido; integración real fuera del alcance \\
  \addlinespace
  Entrega & Integración continua de fallo rápido + despliegue continuo + Docker & Omitir el despliegue continuo (TrackBites) & Decidida; destinos abiertos \\
  \addlinespace
  \multicolumn{4}{@{}l@{}}{\emph{* Alternativa analizada en este capítulo; no consta como barajada en la documentación previa del proyecto.}} \\
\end{longtable}
\endgroup
